% P19-E1: bounded foundations for empirical evidence about reasoning.
\chapter{Questions, Measurements and Empirical Evidence}
\label{sec:p19-evidence-methods}

An empirical study of reasoning connects a question about people to a
specified task, recorded responses and an argument about what those responses
support. Each connection needs justification. A correct answer can be produced
by several processes; a statistical association can have several explanations;
and a result in one task and sample need not transfer to another.

This chapter develops an evidence record and a procedure for planning and
interpreting such studies. Its empirical corpus is deliberately bounded:
Frederick's original Cognitive Reflection Test presentation and Sinayev and
Peters' two studies of reflection and calculation
\cite{Frederick2005CRT,SinayevPeters2015Calculation}.
The American Statistical Association's statement supplies the principles used
to interpret reported \(p\)-values \cite[pp.~131--132]{WassersteinLazar2016ASA}.
The studies were selected for their common task family and their differing
measurement arguments, not to estimate the frequency of favorable results in
the literature. The source cutoff is 7 September 2026. This three-document
corpus is not an exhaustive literature review or a current consensus estimate.

The evidence record, planning procedure and proposed follow-up below are
handbook syntheses. The numerical study results retain their source identities.
No participant-level reanalysis, newly collected data or completed follow-up
experiment is reported here.

\section{Specify the claim before choosing a score}
\label{sec:p19-question-and-measure}

\begin{hbdefinition}{def:p19-empirical-question}{Question, task, construct and measure}
An \emph{empirical question} identifies the people, circumstances and response
or relation to be investigated. A \emph{task} is the actual instruction and
stimulus procedure presented to a participant, including response format,
timing, available information and incentives. A \emph{construct} is the
capacity, disposition or process that the investigator intends to study.
An \emph{operational measure} is the specified rule that converts recorded
responses into a quantity or category.

For example, the number of correct answers to three numerical problems is a
measure. A tendency to inspect an initially attractive answer is a proposed
construct. Equating them requires evidence about how the responses arise;
calling a score a reflection test does not supply that evidence.
\end{hbdefinition}

Three questions that sound similar require different evidence. One may ask
whether two scores are associated in a sample, whether a score predicts
responses from new participants, or whether an intervention changes responses.
The first concerns an observed relation. The second requires a specified
prediction procedure assessed on observations not used to fit or select it.
The third requires a design and assumptions that support an intervention
comparison. An association estimated from the same records used to select a
model does not, by itself, answer the other two questions.

The task should expose the relevant distinction. A final numerical answer
alone does not show whether someone first formed an erroneous intuition,
noticed an inconsistency, guessed, remembered a familiar answer or carried out
a calculation directly. Response times, intermediate answers and explanations
can supply additional observations, but instructions to report them can also
change the task. They need their own recording rules and interpretation.
The two selected empirical papers make the distinction between score and
construct consequential \cite[pp.~26--28, 35]{Frederick2005CRT}
\cite[pp.~2--5, 12--14]{SinayevPeters2015Calculation}.

\begin{hbdefinition}{def:p19-benchmark}{Benchmark and observed response}
A \emph{benchmark} specifies a criterion against which a response is compared,
together with the assumptions that make the criterion applicable. The
\emph{observed response} is what the participant actually reported or did.
The criterion and the observation are different parts of the evidence record.
\end{hbdefinition}

An arithmetic problem can have a uniquely correct answer under its stated
equations. A probability judgment can violate a mathematical constraint if
the participant understands the events in the intended way. For example,
\cref{thm:p15-finite-probability-laws} gives
\(\mathbb P(A)+\mathbb P(A^c)=1\) and
\(\mathbb P(A)\leq\mathbb P(B)\) when \(A\subseteq B\).
Applying these constraints to a person's estimates requires the same outcome
space, the same information and genuinely complementary or nested events.
Different time horizons or interpretations can invalidate the intended
comparison before any numerical calculation begins.

A preference for an immediate payment or a safer prospect has no universally
correct label supplied merely by its frequency or its association with a test
score. The finite ordinal representation of
\cref{def:p17-ordinal-representation} preserves a supplied ordering; it does
not establish a person's utility for money, beliefs about payment, borrowing
constraints or acceptable risk. Choosing the prospect with larger expected
money requires an additional decision criterion if it is to count as better
rather than simply as a coded response. Frederick's discussion explicitly
addresses, and takes positions within, this normative dispute
\cite[pp.~38--41]{Frederick2005CRT}. The empirical associations below do not
settle it.

\section{Plan a record that can constrain the conclusion}
\label{sec:p19-design-and-record}

\begin{hbdefinition}{def:p19-population-and-estimand}{Population, sample and target quantity}
The \emph{target population} is the set or class of people to which the intended
claim refers. The \emph{sampling frame} is the set from which recruitment is
actually attempted. The \emph{recruited sample}, the participants providing a
particular response, and the \emph{analysis sample} remaining after exclusions
and missing-data rules need not coincide.

An \emph{estimand} is the quantity the analysis aims to estimate, with its
population, task and comparison fixed. Examples are a response proportion in
a specified population, a difference between two such proportions, or a
coefficient in a specified statistical model. An \emph{estimate} is the value
computed from the available observations by a declared procedure.
\end{hbdefinition}

Random recruitment and random assignment address different problems. Random
recruitment concerns who is observed; random assignment concerns which
condition an observed participant receives. A convenience sample can be
randomly assigned without becoming representative of a wider population.
Conversely, observing a broad population sample does not turn naturally
occurring differences in scores into randomized interventions.

\begin{hbmethod}{def:p19-evidence-record}{Constructing an empirical evidence record}
\begin{enumerate}
\item State the question and the target quantity. Specify whether the intended
      conclusion is descriptive, predictive, causal or a comparison with a
      normative benchmark. Give the benchmark's assumptions separately.
\item Identify the target population, recruitment frame, recruitment dates,
      setting and incentives. Record how many people were approached,
      recruited, measured and analyzed whenever the source supplies those
      counts; mark unavailable counts as unavailable.
\item Preserve the task version, language, presentation order, response format,
      permitted aids and timing. Record prior task exposure if it matters to
      the proposed interpretation. A changed item set is a changed instrument.
\item Define each response code, score, denominator and missing-data rule.
      State how ambiguous answers, omissions, repeated observations and
      inattentive responding are handled. Preserve reasons for exclusions.
\item Describe the comparison design, assignment procedure and unit of
      analysis. Identify repeated responses within people or groups; their
      number is not automatically the number of independent participants.
\item Report the estimate with its units, analysis sample, model or calculation
      rule and reported uncertainty. Give the adjustment variables and any
      transformations or selection steps that change the estimate's meaning.
\item Examine competing explanations and sensitivity to consequential choices.
      Distinguish analyses planned before inspecting the result from analyses
      selected afterward. Report relevant unfavorable or inconclusive results.
\item End with a conclusion whose population, task, comparison and uncertainty
      match the record. State what a replication or transfer study would need
      to preserve and what it would deliberately change.
\end{enumerate}
This is a reusable reporting procedure, not a certification that completing
its fields proves a theory. Its emphasis on design, measurement, context and
full reporting accords with the ASA statement
\cite[pp.~131--132]{WassersteinLazar2016ASA}.
\end{hbmethod}

Several design choices have immediate inferential consequences. A comparison
between pre-existing score groups leaves differences in education, numeracy,
motivation and experience available as explanations. Adjusting for measured
variables changes the statistical comparison, but does not establish that all
relevant differences were measured accurately. Repeating a task can introduce
learning or order effects. Averaging different items across participants can
change the mixture of tasks represented by a nominally common score. An
exclusion rule chosen because it strengthens a desired result changes the
evidence available for evaluating that result.

Planning therefore includes plausible failure cases: an instruction may not
alter the intended process, two instruments may share arithmetic demands,
dropout may differ between conditions, or a score may have too little variation
to distinguish participants. These possibilities should constrain measurement
and reporting before a result is known.

\section{Describe estimates and uncertainty precisely}
\label{sec:p19-estimates-and-uncertainty}

\begin{hbdefinition}{def:p19-response-summaries}{Response proportions and differences}
For a specified group \(g\), let \(n_g>0\) be the number of eligible recorded
responses to one task and let \(k_g\) be the number receiving a specified code.
The observed proportion is
\[
 \widehat p_g\defeq\frac{k_g}{n_g}.
\]
For two such groups, labelled \(0\) and \(1\), define the observed difference
\(\widehat\Delta\defeq\widehat p_1-\widehat p_0\).
The labels identify the groups; they do not imply random assignment.
A difference of \(0.15\) is 15 percentage points. It is not a 15 percent
relative increase. If a denominator is zero, its proportion is undefined.
\end{hbdefinition}

These are descriptions of observed records. Interpreting them as estimates
for unobserved people requires a sampling or statistical model. Published
rounded percentages support correspondingly approximate arithmetic; they do
not license inventing exact success counts. A denominator must belong to the
specific estimate, rather than to the largest sample mentioned in the paper.

\begin{hbdefinition}{def:p19-uncertainty-terms}{Sampling variability and reported uncertainty}
\emph{Sampling variability} is the variation an estimator would exhibit across
repeated datasets under a specified data-generating model or sampling design.
A \emph{standard error} is an estimate of the estimator's standard deviation
under that model or design. It differs from the standard deviation of individual
responses. A confidence interval comes from a procedure with a specified
coverage property under stated assumptions; its confidence level is not, by
itself, a posterior probability assigned to the particular reported interval.

A \(p\)-value reports a tail probability for a specified statistical summary
under a specified model and extremeness rule, commonly including a null
hypothesis. It does not give the probability that the hypothesis is true or
that chance alone caused the observations
\cite[p.~131]{WassersteinLazar2016ASA}.
\end{hbdefinition}

An effect estimate describes magnitude on a stated scale. Its standard error
or interval addresses uncertainty under the analysis assumptions. Neither
quantifies every uncertainty about task interpretation, measurement error,
nonresponse, population transfer or model choice. Report those limitations in
words and investigate them with appropriate additional observations or
analyses. A small reported \(p\)-value does not establish a large effect, a
correct causal explanation or a practically important result; a large one
does not establish absence of an effect. The ASA statement also requires
attention to selection among analyses and full reporting, rather than using
a threshold as the entire conclusion
\cite[principles~3--6, pp.~131--132]{WassersteinLazar2016ASA}.

The cases below retain the uncertainty actually supplied by their sources.
No unreported confidence interval or participant-level reanalysis is invented.
The full derivation and calibration of inferential procedures, including
regression standard errors, sampling weights and interval construction, belong
to their respective statistical treatments; the purpose here is to state what
a reported number means and what else a reader needs to assess it.

\section{Case one: a response difference on an intertemporal task}
\label{sec:p19-frederick-case}

\paragraph{Question and instrument.}
Frederick introduced a three-item Cognitive Reflection Test, abbreviated CRT.
Its numerical problems have designated correct answers and salient erroneous
answers; a total score counts correct answers from zero to three. The proposed
interpretation concerns resistance to an initially attractive response, while
the observed variable is correctness on those items. The original paper itself
recognizes contributions from mathematical skill and reading comprehension
\cite[pp.~26--28, 35]{Frederick2005CRT}.

\paragraph{Population, recruitment and comparison.}
The paper reports 3,428 respondents across 35 studies during 26 months beginning
in January 2003. Most were undergraduates at universities in the midwestern
and northeastern United States, with other recruitment settings listed in
Table~1. The usual procedure paid eight dollars for a 45-minute questionnaire;
the paper identifies exceptions for class participation, a fireworks gathering
and web recruitment. These are not a probability sample of all adults.
The primary comparison in many tables is between respondents scoring zero and
those scoring three; it omits the intermediate score groups from that contrast
\cite[pp.~28--29 and notes to table~1]{Frederick2005CRT}.

\paragraph{Task, response and estimate.}
Table~2, item~b, pairs an immediate payment of 100 dollars with 140 dollars
after one year. It is one of the hypothetical monetary choices identified in
the source. The response coded as patient is selection of the later payment.
The source reports 22 percent for the zero-score group, with 409 respondents,
and 37 percent for the three-score group, with 297 respondents. Thus the
difference calculated from the reported percentages is
\[
          \widehat\Delta\approx0.37-0.22=0.15.
\]
The approximately 15 percentage point contrast concerns the 706 respondents
in those two groups who answered that item, not all 3,428 respondents or all
35 studies separately. The source reports \(p<0.0001\) for the group
comparison and specifies chi-square tests for dichotomous items. The printed
percentages are rounded; the table does not supply an exact confidence interval
for this difference \cite[pp.~30--31, table~2b]{Frederick2005CRT}.

\paragraph{What the result establishes.}
In these observed groups, reporting all three correct CRT answers was associated
with more frequent selection of this later hypothetical payment. The groups
were formed from test responses, not by random assignment to a reflective
state. The contrast does not estimate the effect of teaching reflection, and
does not show that everyone with a higher score is more patient. Numeracy,
education, experience, interpretation of the question and other differences
can contribute to the relation.

\paragraph{Uncertainty and rival interpretations.}
The reported test addresses the stated statistical comparison. Selection into
the studies, different questionnaire contexts, exclusions and item-specific
nonresponse require additional information before a population estimate or a
new confidence interval could be justified. The paper discusses reasons that
choices between dated payments can reflect beliefs about receipt, changing
wealth, inflation or opportunity costs as well as valuation of delay. Those
considerations are hypotheses about interpretation, not measured explanations
for every response in this item \cite[pp.~31--32]{Frederick2005CRT}.

Selecting this short-horizon item also must not conceal the source's limiting
results: several other intertemporal items have weaker or nonsignificant
reported group differences, including nonmonetary and longer-horizon cases.
The one-item contrast is an illustrative evidence record, not an estimate of
an invariant effect across all intertemporal decisions
\cite[table~2 and discussion, pp.~30--32]{Frederick2005CRT}.

\paragraph{Normative and transfer boundary.}
More frequent selection of the later payment is the observed outcome. Calling
it the rational response would require additional assumptions about the
participant's alternatives, resources, beliefs and objectives. Replication
would need to specify whether it preserves the amounts, delay, hypothetical
status, wording, score groups and recruitment setting. A result with real
payments or new task items would answer a related, explicitly changed question.

\section{Case two: what does the score measure?}
\label{sec:p19-measurement-case}

Sinayev and Peters examined whether associations attributed to reflection could
instead be accounted for by numerical abilities. Their response coding
distinguishes a designated intuitive error, a correct answer, and another
error. Their psychological interpretation posits an initial opportunity to
avoid the intuitive response followed by calculation. The categories are
observable codes; that sequence of mental operations is a model
\cite[pp.~3--5, table~1]{SinayevPeters2015Calculation}.

\begin{hbdefinition}{def:p19-response-decomposition}{Two scores from three response categories}
For one participant's \(m>0\) completed and classifiable items, let
\(n_I,n_C,n_O\) count designated intuitive errors, correct answers and other
errors, respectively, with \(n_I+n_C+n_O=m\). Define
\[
 \widehat a\defeq\frac{n_C+n_O}{m},\qquad
 \widehat c\defeq\frac{n_C}{n_C+n_O}
                       \quad\text{when }n_C+n_O>0.
\]
The first is the proportion of answers that avoid the designated intuitive
error. The second is the proportion correct among those answers. They are
the proportion-based measures called Cognitive Reflection and Calculation
in the source. These descriptive names do not prove a unique psychological
interpretation of either score. If the second denominator is positive,
direct cancellation gives \(\widehat a\widehat c=n_C/m\).
\end{hbdefinition}

Guessing or inattentive responses can increase \(\widehat a\) without
demonstrating successful reflection. Knowing an answer beforehand can produce
a correct response without the posited sequence. Conversely, avoiding the
designated error does not guarantee correct calculation. For the original
two-price problem, Frederick describes correct answers as nearly universal
among those avoiding the designated error. Sinayev and Peters' Study~1
Table~2 reports 78 percent designated errors, 14 percent correct answers and
8 percent other errors on that item. From these rounded proportions, correctness
among the non-designated answers is approximately
\[
                    \frac{14}{14+8}\approx0.64.
\]
This is a calculation from one published item's aggregate proportions, not
the mean of participants' \(\widehat c\) scores. The substantial other-error
category in this sample limits the claim that detecting the salient error is
effectively sufficient for solving the problem. Differences in samples and
contexts prevent this contrast from identifying a single cause of the
discrepancy \cite[p.~27]{Frederick2005CRT}
\cite[pp.~5--7, table~2]{SinayevPeters2015Calculation}.

\paragraph{A complete Study 1 evidence record.}
Study~1 collected internet responses from 1,413 participants in the Understanding
America Study between 31 May and 22 October 2014, using address-based recruitment.
Participants received ten dollars for a survey taking about half an hour.
The median participant was 49, had an associate degree and reported income
between 50,000 and 60,000 dollars; 52 percent were female. Four participants
were explicitly excluded for response patterns interpreted as inattention.
The complete-case regression analysis used 1,225 participants completing the
relevant tasks. These are different denominators; the difference between the
initial and complete-case counts must not all be called an attention exclusion
\cite[pp.~4--7]{SinayevPeters2015Calculation}.

The instrument had five CRT items: two from the original test and three
additional items. Numeracy was scored from six non-CRT items in an eight-item
scale, then centered and standardized. The Study~1 regression predictors for
reflection and calculation come from the following response model, rather than
from the raw CRT total or the proportion scores above
\cite[pp.~4--7]{SinayevPeters2015Calculation}.

\begin{hbdefinition}{def:p19-study1-response-model}{Study 1 response model and fitted person scores}
Index participants by \(i\) and the five items by \(j\in\{1,\ldots,5\}\).
Write \(X_{ij}\in\{\mathrm I,\mathrm C,\mathrm O\}\) for a classified
response: designated intuitive error, correct answer, or other error.
The source's first binary step records whether the response avoids
\(\mathrm I\). Its second records correctness conditional on avoiding
\(\mathrm I\). An \(\mathrm I\) response is coded zero at the first step
and missing at the second; \(\mathrm C\) is coded \((1,1)\), and
\(\mathrm O\) is coded \((1,0)\). Missing here represents a step that the
model assumes was not reached, rather than a second observed failure.

In notation for the source's item-and-person log-odds model, let \(r_i,c_i\)
be the real person effects for the first and second step. Let
\(\alpha_j,\beta_j\) be their respective item intercepts and put
\(L(s)\defeq1/(1+\exp(-s))\) for real \(s\), the logistic function.
Superscripts \(R,C\) distinguish reflection and calculation, and \(\Pr\)
denotes probabilities in this fitted response model. The step probabilities,
conditional on the person effects, are
\[
 \pi^{R}_{ij}\defeq L(\alpha_j+r_i),\qquad
 \pi^{C}_{ij}\defeq L(\beta_j+c_i).
\]
The second probability is conditional on reaching the calculation step.
Consequently the three response probabilities are
\[
\begin{aligned}
 \Pr(X_{ij}=\mathrm I\mid r_i,c_i)&=1-\pi^{R}_{ij},\\
 \Pr(X_{ij}=\mathrm C\mid r_i,c_i)&=\pi^{R}_{ij}\pi^{C}_{ij},\\
 \Pr(X_{ij}=\mathrm O\mid r_i,c_i)&=\pi^{R}_{ij}(1-\pi^{C}_{ij}).
\end{aligned}
\]
The ten intercepts are the five item effects at each step, called difficulty
parameters in the source.
A larger intercept increases the corresponding success probability.
The same \(r_i\) enters all five first steps and the same \(c_i\) all five
second steps. Model~2 estimates separate variation in these two person effects
and their correlation; Model~1 restricts the person variation to one common
effect across both steps. The source reports Model~2 person-effect scales
1.2 for reflection and 2.1 for calculation, with correlation 0.40
\cite[pp.~5--7, tables~1 and 4]{SinayevPeters2015Calculation}.

After fitting Model~2, the authors obtain each participant's reflection and
calculation predictors from the conditional distributions of the person
effects given that participant's responses and the fitted model. The source
identifies the extracted values as conditional modes, meaning the most likely
person-effect values under that fitted model. Denote these reported scores
by \(\widehat r_i,\widehat c_i^{\rm mod}\). The modal calculation score
is different from the proportion \(\widehat c\) defined earlier; neither
modal score is a response probability or a count of correct answers.

These scores use the model's log-odds deviation scale relative to the fitted
item intercepts. A one-unit increase in a person effect multiplies the odds
of success at its corresponding step by \(\exp(1)\), with the item fixed;
it does not multiply that success probability by \(\exp(1)\). The paper
reports no further standardization of the extracted modal predictors to unit
variance. Table~3 gives their observed means and standard deviations as
\(-0.01,0.86\) for reflection and \(0.04,1.27\) for calculation,
respectively. Those empirical summaries of extracted scores are distinct from
the model's estimated person-effect scales 1.2 and 2.1. Numeracy, in contrast,
is explicitly centered and standardized and has mean zero and standard
deviation one in that table \cite[pp.~4, 6--7]{SinayevPeters2015Calculation}.

The paper names the \texttt{nlme} package, version 3.1, and gives the response
coding, item/person structure, log-odds interpretation and conditional-mode
description. It does not print its fitting or score-extraction calls, numerical
integration and optimization settings, or participant score vectors. It also
does not write out the random-effects density and all identification constraints
in its own specification, instead referring to the earlier response-model
literature. Thus the reported predictor construction and units are identified,
but the exact numerical scores cannot be reconstructed from these summary
tables alone. In particular, the source does not specify whether its modal
extraction is implemented jointly or marginally; no such choice is supplied
here as an invented estimator. The displayed probability equations express
the stated two-step log-odds structure, not an independently rerun fit.
\end{hbdefinition}

One outcome counted inconsistencies in four pairs of complementary probability
judgments. Each pair was scored inconsistent when its two percentage estimates
summed to a value at least ten percentage points away from 100. The total
therefore ranges from zero to four. The mathematical benchmark is the
complement rule, conditional on the intended common event interpretation.
The ten-point tolerance is an operational scoring choice, not a modification
of probability theory. The authors report also examining tolerances of five,
fifteen and twenty points, and retaining ten because it gave the most favorable
test of their reflection hypothesis. This disclosed choice belongs in the
analysis history; it must not be recast as a verified advance commitment
\cite[p.~4]{SinayevPeters2015Calculation}.

For an explicitly hypothetical handbook respondent, complementary estimates
of 65 and 45 percent sum to 110 percent and meet that inconsistency threshold.
Estimates of 65 and 40 percent sum to 105 percent and do not meet it, although
they still fail the exact complement equation. A tolerance-based response code
and the mathematical benchmark therefore answer different questions. These
two pairs are illustrations of the scoring rule, not observations from the study.

The unadjusted correlation between the model-derived reflection score and this
inconsistency count is \(-0.20\). Its Table~3 cell does not state a separate
sample size; the regression sample must not silently be assigned to that cell.
Table~5 reports the following unstandardized
regression coefficients, with standard errors in parentheses. Its reflection
and calculation rows use \(\widehat r_i\) and
\(\widehat c_i^{\rm mod}\), respectively. Both columns
also include age, income, education and gender; the second additionally
includes numeracy \cite[tables~3 and 5, pp.~6--8]{SinayevPeters2015Calculation}.
\[
\begin{array}{l|rr}
 \text{Predictor}&\text{Without numeracy}&\text{With numeracy}\\\hline
 \text{Reflection score}&-0.07\ (0.05)&-0.06\ (0.05)\\
 \text{Calculation score}&-0.12\ (0.03)&-0.09\ (0.03)\\
 \text{Numeracy}&\text{not included}&-0.13\ (0.04)
\end{array}
\]
The source marks the calculation coefficients and the included numeracy
coefficient as significant at its \(0.05\) threshold; it does not mark either
reflection coefficient. This pattern concerns conditional associations in
these fitted models. It does not demonstrate that reflection never matters.

The outcome is a count, so the numeracy coefficient of \(-0.13\) describes
0.13 fewer inconsistent pairs per one-standard-deviation increase in the
standardized numeracy predictor, conditional on the other predictors in this
linear model. It is not a 13 percentage point decrease in every participant's
error probability. The reflection and calculation scores have their own
model-derived scales; comparing the absolute coefficients alone would not
compare effects on a common scale. The reported standard errors quantify
uncertainty under the fitted analysis, not uncertainty about the existence or
causal separateness of the proposed mental processes.

\paragraph{Model comparison, robustness and warranted conclusion.}
The two-component response model fitted better than the specified one-component
alternative. The Bayesian information criterion, abbreviated BIC, compares
model fit with a penalty for model complexity. The paper reports a BIC
difference of 198 and a chi-square comparison statistic of 217 with two degrees
of freedom, \(p<0.001\). These reported comparisons concern those
models and their assumptions. They do not exclude unconsidered models, prove
that every participant first generated an intuition, or turn the fitted
components into experimentally isolated mechanisms
\cite[p.~5]{SinayevPeters2015Calculation}.

The appendix's robust-regression version of the adjusted inconsistency model
reports reflection \(-0.05\ (0.05)\), calculation \(-0.10\ (0.04)\), and
numeracy \(-0.11\ (0.04)\), retaining the broad pattern. This is a sensitivity
analysis of the study data, not an independent replication. Its agreement
addresses one possible sensitivity to influential observations; it does not
resolve selection, construct validity or every missing-data assumption
\cite[appendix, table~A1, p.~16]{SinayevPeters2015Calculation}.

The warranted conclusion is that this study's CRT-derived reflection measure
did not contribute a clearly detected independent association with the selected
inconsistency outcome after the stated adjustments, whereas numerical measures
did. The authors' account challenges a uniquely reflective explanation of CRT
associations. The study does not identify the causal effect of increasing
numeracy, nor prove that the observed response code is a pure measure of
reflection. Competing explanations include shared task demands, imperfectly
measured abilities, correlated predictors, task interpretation, and the
selection of complete cases. The reported associations and robustness check
constrain these possibilities without eliminating them all.

\section{Compare studies before combining their results}
\label{sec:p19-comparability-and-replication}

\begin{hbwarning}{def:p19-composite-comparability}{A common score range is not a common measurement}
Frederick's decision composites average the available coded responses for a
participant, while the items and accompanying cognitive tests differ across
studies. His footnote~12 explicitly explains why this complicates comparisons
between cognitive measures and describes comparisons restricted to respondents
who received both measures \cite[p.~36]{Frederick2005CRT}.

Sinayev and Peters' Study~2 similarly permits a composite to average different
available components for different participants. Its decision composite combines
standardized probability-consistency, confidence and choice measures; it is not
a simple percentage of universally incorrect decisions
\cite[pp.~8--10]{SinayevPeters2015Calculation}.
Identical numerical ranges or a shared composite name do not establish that
two people, samples or instruments were measured in the same way.
\end{hbwarning}

A comparability assessment should align the actual item set, response coding,
incentives, population, analysis sample and estimand. The original CRT contains
three numerical items. The follow-up's Study~1 uses five items and derives
scores from a latent response model; Study~2 uses three items and proportion
scores. Study~1 has four subset/superset comparisons and eight time comparisons,
whereas Study~2 has six and ten, respectively; both have four complementary
pairs. Study~2 also includes incentivized choices and draws from RAND's American
Life Panel. The paper dates that study's data to 2006 in its discussion, despite
its publication after the 2014 data of Study~1
\cite[pp.~4, 8--10, 13]{SinayevPeters2015Calculation}.
Publication order and data-collection order therefore must be recorded separately.

For Study~2's complementary-pair outcome, the paper reports an unadjusted
reflection correlation of \(-0.12\), followed by an adjusted reflection
coefficient of \(-0.01\) with \(p=0.81\); Table~7 gives 926 observations for
that outcome. This selected model adjusts for Calculation, numeracy, age,
income, education and gender. It does not include the intelligence composite
used in the paper's later analyses of composite outcomes. This is a repeat of
the broad association pattern in another
sample under the specified measurements. The changed predictor construction
means that its slope should not be numerically pooled with the Study~1 latent
score slope as though both estimated an identical quantity
\cite[pp.~10--11]{SinayevPeters2015Calculation}.

\begin{hbwarning}{def:p19-undefined-calculation-score}{An unobserved conditional score is not a zero probability}
If all of a participant's completed CRT responses are designated intuitive
errors, then \(n_C+n_O=0\), so \(\widehat c\) is undefined. The data contain
no non-designated responses from which to estimate correctness conditional on
such a response. The issue is the same positive-denominator requirement made
explicit in \cref{def:p15-null-conditioning}.

In Study~2, about 24 percent of participants had this response pattern. The
authors assigned a calculation score of zero and reported essentially identical
results under an alternative imputation using other measured variables
\cite[p.~9 and notes~2--3]{SinayevPeters2015Calculation}.
The zero is a declared scoring convention, not the value of \(0/0\) and not
an observed conditional success rate. A reuse of the data must retain that
distinction and report sensitivity to the convention.
\end{hbwarning}

\paragraph{Reliability, familiarity and competing accounts.}
The follow-up discusses low internal consistency of several outcomes and the
possibility that the instruments measure several constructs. Internal
consistency concerns relations among item responses; it does not establish
that the score measures the intended psychological process. The paper also
raises prior exposure as a possible problem for the familiar CRT, while arguing
that exposure was unlikely in its own samples because of its item choices and
collection dates. That argument is not a direct experimental estimate of an
exposure effect in these samples \cite[pp.~7--9, 12--14]{SinayevPeters2015Calculation}.

These two empirical articles cannot settle all competing accounts of reflection,
numeracy, memory and motivation. A new score or a better fitting model may
improve an explanation, but a change in measured association alone does not
identify which mental process changed. The preserved distinctions support
targeted follow-up questions instead of a count of nominally positive and
negative studies.

\section{A proposed follow-up with an explicit inferential target}
\label{sec:p19-proposed-design}

\begin{hbexample}{def:p19-proposed-factorial-design}{Separating an assignment effect from a process claim}
The following is a hypothetical handbook design, not a study that has been
conducted. Its primary question is whether adding an instruction to check a
first response reduces complementary-probability inconsistencies when a
numerical aid is already supplied, on new tasks in a specified recruited adult
sample. A separate secondary question concerns the checking effect averaged
equally over the aid-present and aid-absent conditions.

Recruit adults able to read the task language from an explicitly recorded
frame. Obtain task-exposure history and the chosen baseline numeracy measure
before assignment, using the same sequence for everyone. Randomly assign each
participant to one of four conditions crossing two factors: a checking
instruction or its matched neutral instruction, and a numerical aid or its
matched control. Fix the content, presentation time, incentives and response
format in advance. Use new items that preserve the intended logical relations,
with a pilot checking whether participants understand those relations.

The primary outcome is each participant's number of inconsistent complementary
pairs on the same four tasks, under a fixed tolerance rule. Preserve all
responses and omission codes. Secondary outcomes, such as response time and
accuracy on an independent arithmetic task, have separate definitions and
analysis roles. Fix the recruitment stopping rule and sample-size justification
using the smallest effect worth distinguishing and an explicit variance and
missingness model; this example supplies no fabricated power calculation.

Let \(\overline Y_{ab}\) denote the observed mean primary outcome in condition
\((a,b)\), where \(a=1\) indicates the checking instruction, \(a=0\) its
control, \(b=1\) the numerical aid and \(b=0\) its control. If all four means
are available, both comparisons below are defined. The primary aid-present
checking contrast requires only that its two constituent means exist:
\[
 \widehat\Delta_{\rm aid}\defeq\overline Y_{11}-\overline Y_{01}.
\]
The retained secondary equal-weight contrast averages checking differences
across the two aid conditions:
\[
 \widehat\Delta_{\rm check}\defeq
 \frac{(\overline Y_{10}-\overline Y_{00})+
       (\overline Y_{11}-\overline Y_{01})}{2}.
\]
For each specified comparison, a negative difference means fewer scored
inconsistencies under the checking instruction. The two contrasts need not
even have the same sign. For hypothetical means
\((\overline Y_{00},\overline Y_{10},\overline Y_{01},\overline Y_{11})
=(4,0,0,2)\), the aid-present difference is \(2-0=2\), while the
equal-weight average is \(((0-4)+(2-0))/2=-1\). Checking worsens the
aid-present outcome in this example despite reducing the equal-weight average.
These are declared hypothetical means, not study observations. Report the
two constituent differences as well as the average, since the aid could change
the instruction's effect. Estimate uncertainty using the actual assignment and
analysis procedure, accounting for participant-level clustering of task
responses. Do not treat four answers from one person as four independently
assigned participants.

Under the specified random assignment, adherence to the assigned conditions,
stable versions of each intervention, no relevant interference between
participants and an adequate treatment of missing outcomes, the design can
support a comparison of the effects of assignment in the recruited sample.
If people do not follow instructions, an analysis retaining assignment groups
concerns the effect of assignment, not automatically the effect of compliance.
If outcome loss differs by condition, report that loss and examine how it could
change the comparison before claiming a treatment effect.

Even a detected assignment effect would not by itself prove that reflection
caused the improvement: the instruction may change effort, caution, time spent
or interpretation. Process measures, alternative control instructions and
predictions that distinguish these accounts would supply further evidence.
Generalization to other ages, languages, tasks or everyday decisions requires
additional sampling and transfer studies. The design states a concrete
intervention question while keeping those broader claims separate.
\end{hbexample}

\section{Write the conclusion the evidence supports}
\label{sec:p19-warranted-conclusion}

A warranted conclusion names the measured relation or assignment comparison,
its sample and task, the estimate and available uncertainty, and the explanations
still compatible with the design. In the first case, the observed CRT groups
differ in one hypothetical payment choice. In the second, the selected
model-derived reflection score did not supply a clearly detected independent
association with the selected inconsistency count after the stated adjustments. Neither result
is a universal ranking of people by rationality or an identified causal effect
of a cognitive ability.

To make a study reproducible, preserve the task version, coding rules,
exclusions, source data and analysis instructions to the extent permitted by
the data's terms. To evaluate replication, record which features were preserved
and which were changed before comparing results. A reanalysis of the same
records, a new sample using the same procedure and a transfer study using new
tasks answer different questions. Agreement strengthens only the claims their
shared assumptions support; disagreement calls for inspecting measurement,
sampling, design and uncertainty before choosing a favored explanation.

This unit supplies a first evidence-method foundation and bounded treatments
from three sources, not the complete methods of empirical reasoning research.
Formal sampling and causal-identification theory, measurement models,
uncertainty derivations, systematic review and meta-analysis, registered
replication, cross-cultural and developmental transfer, and the remaining P19
subjects require their own complete treatments. The ASA statement itself leaves
many statistical controversies and procedures outside its stated scope
\cite[pp.~130--132]{WassersteinLazar2016ASA}.
